\section{Probabilistic Modelling and Inference}\label{sec:modelling}
\subsection{Bayesian Logistic Regression}\label{sec:first_model}
I use a series of increasingly complex logistic regression models, each of the following form:
\[
\theta_i
=
\alpha
+ \mathbf{\alpha}_{\mathrm{cat}}(\mathbf{c}_i)
+ \boldsymbol{\beta}^\top \mathbf{x}_i,
\qquad
y_i \sim \mathrm{Bernoulli}\!\left(\mathrm{logistic}(\theta_i)\right)
\]
with priors
\[
\alpha_{\mathrm{cat}}(j) \sim \mathcal{N}(0,1) \quad \forall \, j\in \mathrm{cat}, \qquad 
\boldsymbol{\beta} \sim \mathcal{N}(0,1), \qquad
\alpha \sim \mathcal{N}(0,2)
\]
where \(\mathbf{x}_i\) denotes continuous covariates and \(\mathbf{\alpha}_{\mathrm{cat}}(\mathbf{c}_i)\) denotes the contribution of categorical variables (e.g.\ province or protection zone), $j$ is a value for a given category $\mathrm{cat}$, and $\alpha$ is the intercept. Priors were chosen to have full support with slightly more variation for $\alpha$ to allow for a low baseline evacuation probability. A hierarchical model could have been employed but I found this increased computational complexity while $N=15,203$ was sufficient to estimate separate category effects directly. Below, I only state the changes relative to the previous for simplicity.\\
\textbf{Minimal:} only fire size and whether the fire is on a FN reserve.
\[
\mathbf{x}_i = \big(\log \text{fire\_size},\ \text{fn\_indicator}_i\big),
\qquad
\mathbf{\alpha}_{\mathrm{cat}}(\mathbf{c}_i) = 0
\]

\textbf{Provinces:} add provinces to the categorical variables.
\[
\mathbf{\alpha}_{\mathrm{cat}}(\mathbf{c}_i) = \alpha_{\mathrm{prov}[i]}
\]

\textbf{Full:} add distance to nearest FN reserve and the protection zone variables.
\[
\mathbf{x}_i = \big(\log \text{fire\_size},\ \text{log\_dist\_to\_fn}_i,\ \text{fn\_indicator}_i\big),
\qquad
\mathbf{\alpha}_{\mathrm{cat}}(\mathbf{c}_i)
=
\alpha_{\mathrm{prov}[i]} + \alpha_{\mathrm{prot}[i]}
\]

\textbf{Responses:} add the response type to the categorical variables.
\[
\mathbf{\alpha}_{\mathrm{cat}}(\mathbf{c}_i)
=
\alpha_{\mathrm{prov}[i]} + \alpha_{\mathrm{prot}[i]} + \alpha_{\mathrm{resp}[i]}
\]

\textbf{Spatial (VI):} remove response type and add region. This model was too computationally expensive to run inference using MCMC thus I used variational inference (VI) to approximate the posterior using the default cmdstanpy parameters ('meanfield' algorithm, independence assumption).

\[
\mathbf{x}_i = \big(\log(\text{fire\_size}_i),\ \text{log\_dist\_to\_fn}_i,\ \text{fn\_indicator}_i\big),
\qquad
\mathbf{\alpha}_{\mathrm{cat}}(\mathbf{c}_i)
=
\alpha_{\mathrm{prov}[i]} + \alpha_{\mathrm{prot}[i]} + u_{\mathrm{region}[i]}.
\]
with a spatial prior imposed on \(u_{\mathrm{region}}\). Let $E$ denotes the set of neighbouring regions. The prior imposes the belief that neighbouring regions have related rates of evacuations.
\[
u_{\mathrm{region}} \sim \mathcal{N}(0,1), \qquad
\sum_{r=1}^R u_r = 0, \qquad
p(\mathbf{u}) \propto
\exp\left(
-\frac{1}{2\sigma_{\mathrm{region}}^2}
\sum_{(r,s)\in E} (u_r - u_s)^2
\right),
\]
\subsection{Diagnostics}
I inspected trace and rank plots for all MCMC models. An example for the Full model with parameters $\alpha$ and $\beta_{\mathrm{fn}}$ is shown in Appendix~\ref{sec:trace_and_rank_plots}, Figure~\ref{fig:full_trace_rank}. I also ensured that $\hat{R}\approx 1$ for all parameter values indicating that all chains agree, and found that the effective sample size (ESS) was over 700 for all parameters on all models in both tail and bulk, with most values between 1000 and 3000, indicating good convergence. I ran posterior predictive checks (PPC) for each model, verifying that a credible interval (CI) over the total number of evacuations is consistent with the data. These can be seen below in Table~\ref{tab:model_comparison_loo}. Finally, I ran a brief sensitivity analysis on the variances of the priors for the Province model in Appendix~\ref{sec:sensitivity_analysis} which indicated that parameter means were stable enough across priors to leave the conclusions unchanged.

\subsection{Model Comparison}\label{sec:model_comparison}
I compare the nested logistic models using Leave-one-out (LOO) expected log predictive density (ELPD). That is, $\text{ELPD} = \sum_{i=1}^{N} \log p(y_i | y_{-i})$. Higher ELPD indicates better out-of-sample predictive performance. Table~\ref{tab:model_comparison_loo} reports ELPD with its standard error (SE), the ELPD difference from the best model, the standard error of that difference (DSE) and a 95\% quantile-based CI for total number of evacs. All Pareto-k values used for estimation of the ELPD were in the 'good' category ($<0.7$) indicating stability of the estimation.

\begin{table}[h!]
\centering
\caption{Model comparison using LOO ELPD.}
\label{tab:model_comparison_loo}
\begin{tabular}{lrrrr}
\hline
Model & ELPD $\pm$ SE & ELPD diff & DSE & Mean total evacs [95\% CI]  \\
\hline
Response & \textbf{-1865} $\pm$ 63.0 & 0.0 & 0.0 & 532 [475, 593] \\
Spatial (VI) & -1868 $\pm$ 63.3 & 2.8 & 9.2 & 538 [448, 636] \\
Full & -1871 $\pm$ 63.1 & 5.6 & 3.5 & 532 [473, 592] \\
Province & -2019 $\pm$ 66.5 & 154.0 & 17.6 & 532 [471, 595] \\
Minimal & -2247 $\pm$ 73.8 & 381.8 & 28.8 & 532 [472, 595] \\
\hline
\end{tabular}
\end{table}

All models posterior CIs for the total number of evacuations are sensible, with the true value at $531$. Note the Spatial model with VI has a wider CI interval and mean slightly farther than the rest from the true value. Response gives the best ELPD, while Spatial (VI) and Full are close with differences only 1–2 SEs apart. There is a large and meaningful gap between Province and the top three models, and an even larger one between Minimal and the rest.

% More detailed analysis of coefficients can be found in Section~\ref{sec:results}.
